\chapter{2007年安徽卷（理）}

\section{单选题}

\begin{problem}
下列函数中，反函数是其自身的函数为（\quad）

\choices
    {\(f(x) = x^{2}\)（\(x \in [0, +\infty)\)）}
    {\(f(x) = x^{3}\)（\(x \in (-\infty, +\infty)\)）}
    {\(f(x) = \e^{x}\)（\(x \in (-\infty, +\infty)\)）}
    {\(f(x) = \frac{1}{x}\)（\(x \in (0, +\infty)\)）}
\end{problem}

\begin{answer}
D
\end{answer}

\begin{solution}
对于选项 A，函数在 \([0,+\infty)\) 上的反函数为 \(f^{-1}(x)=\sqrt{x}\)，不是原函数；对于选项 B，函数 \(f(x)=x^3\) 的反函数为 \(f^{-1}(x)=\sqrt[3]{x}\)，不是原函数；对于选项 C，函数 \(f(x)=\e^x\) 的反函数为 \(f^{-1}(x)=\ln x\)，不是原函数. 对于选项 D，若 \(y=\frac{1}{x}\)，则交换 \(x\)，\(y\) 后仍得到 \(y=\frac{1}{x}\)，所以其反函数就是自身. 故选 D.
\end{solution}

\begin{problem}
设 \(l\)，\(m\)，\(n\) 均为直线，其中 \(m\)，\(n\) 在平面 \(\alpha\) 内，“\(l \perp \alpha\)”是“\(l \perp m\) 且 \(l \perp n\)”的（\quad）

\choices
    {充分不必要条件}
    {必要不充分条件}
    {充分必要条件}
    {既不充分也不必要条件}
\end{problem}

\begin{answer}
A
\end{answer}

\begin{solution}
若 \(l\perp\alpha\)，则由直线与平面垂直的判定性质，\(l\) 垂直于平面 \(\alpha\) 内的任意直线，因而 \(l\perp m\) 且 \(l\perp n\). 所以“\(l\perp\alpha\)”是“\(l\perp m\) 且 \(l\perp n\)”的充分条件. 反之，平面 \(\alpha\) 内的两条直线不一定相交，例如取 \(m\)，\(n\) 为平面内两条平行直线，另取平面内与它们垂直的直线 \(l\)，则 \(l\perp m\) 且 \(l\perp n\)，但 \(l\) 不垂直于平面 \(\alpha\). 因此该条件不是必要条件. 故选 A.
\end{solution}

\begin{problem}
若对任意 \(x \in \R\)，不等式 \(|x| \geqslant \alpha x\) 恒成立，则实数 \(\alpha\) 的取值范围是（\quad）

\choices
    {\(\alpha < -1\)}
    {\(|\alpha| \leqslant 1\)}
    {\(|\alpha| < 1\)}
    {\(\alpha \geqslant 1\)}
\end{problem}

\begin{answer}
B
\end{answer}

\begin{solution}
当 \(x>0\) 时，\(|x|\geqslant\alpha x\) 化为 \(x\geqslant\alpha x\)，所以 \(\alpha\leqslant1\). 当 \(x<0\) 时，\(|x|=-x\)，不等式为 \(-x\geqslant\alpha x\)，除以负数 \(x\) 后不等号改变方向，得 \(\alpha\geqslant-1\). 当 \(x=0\) 时不等式恒成立. 因此 \(-1\leqslant\alpha\leqslant1\)，即 \(|\alpha|\leqslant1\). 故选 B.
\end{solution}

\begin{problem}
若 \(a\) 为实数，\(\frac{2 + a\i}{1 + \sqrt{2}\i} = -\sqrt{2}\i\)，则 \(a\) 等于（\quad）

\choices
    {\(\sqrt{2}\)}
    {\(-\sqrt{2}\)}
    {\(2 \sqrt{2}\)}
    {\(-2\sqrt{2}\)}
\end{problem}

\begin{answer}
B
\end{answer}

\begin{solution}
由题意，
\[
2+a\i=(-\sqrt{2}\i)(1+\sqrt{2}\i)=2-\sqrt{2}\i.
\]
比较实部和虚部，得到 \(a=-\sqrt{2}\). 故选 B.
\end{solution}

\begin{problem}
若 \(A = \{x \in \Z \mid 2 \leqslant 2^{2 - x} < 8\}\)，\(B = \{x \in \R \mid |\log_2 x| > 1\}\)，则 \(A \cap (\complement_{\R} B)\) 的元素个数为（\quad）

\choices
    {0}
    {1}
    {2}
    {3}
\end{problem}

\begin{answer}
C
\end{answer}

\begin{solution}
由 \(2\leqslant2^{2-x}<8=2^3\)，且指数函数 \(2^x\) 递增，得 \(1\leqslant2-x<3\)，即 \(-1<x\leqslant1\). 由于 \(x\in\Z\)，所以 \(A=\{0,1\}\).

由 \(|\log_2x|>1\)，得 \(\log_2x>1\) 或 \(\log_2x<-1\)，从而 \(B=(0,\frac12)\cup(2,+\infty)\). 因此 \(\complement_{\R}B=(-\infty,0]\cup[\frac12,2]\)，故 \(A\cap(\complement_{\R}B)=\{0,1\}\)，其元素个数为 \(2\). 故选 C.
\end{solution}

\begin{problem}
函数 \( f(x) = 3\sin \left(2x - \frac{\pi}{3}\right) \) 的图象为 \(C\)，

\begin{circlelist}
\item 图象 \(C\) 关于直线 \( x = \frac{11\pi}{12} \) 对称；
\item 函数 \( f(x) \) 在区间 \( \left(-\frac{\pi}{12}, \frac{5\pi}{12}\right) \) 内是增函数；
\item 由 \(y = 3\sin 2x\) 的图象向右平移 \(\frac{\pi}{3}\) 个单位长度可以得到图象 \(C\).
\end{circlelist}

以上三个论断中，正确论断的个数是（\quad）

\choices
    {0}
    {1}
    {2}
    {3}
\end{problem}

\begin{answer}
C
\end{answer}

\begin{solution}
对于论断 \(\circled{1}\)，函数图象的对称轴满足
\[
2x-\frac{\pi}{3}=\frac{\pi}{2}+k\pi\quad(k\in\Z),
\]
所以 \(x=\frac{5\pi}{12}+\frac{k\pi}{2}\). 取 \(k=1\)，得到 \(x=\frac{11\pi}{12}\)，论断 \(\circled{1}\) 正确.

函数的导数为 \(f'(x)=6\cos(2x-\frac{\pi}{3})\). 当 \(-\frac{\pi}{12}<x<\frac{5\pi}{12}\) 时，\(-\frac{\pi}{2}<2x-\frac{\pi}{3}<\frac{\pi}{2}\)，所以 \(f'(x)>0\)，论断 \(\circled{2}\) 正确.

将 \(y=3\sin2x\) 的图象向右平移 \(\frac{\pi}{3}\) 个单位得到 \(y=3\sin(2x-\frac{2\pi}{3})\)，而 \(f(x)=3\sin(2x-\frac{\pi}{3})\). 要得到图象 \(C\)，应向右平移 \(\frac{\pi}{6}\) 个单位，故论断 \(\circled{3}\) 错误. 因此正确论断有两个. 故选 C.
\end{solution}

\begin{problem}
如果点 \(P\) 在平面区域 \(\begin{cases}
2x - y + 2 \geqslant 0,\\
x - 2y + 1 \leqslant 0,\\
x + y - 2 \leqslant 0
\end{cases}\) 上，点 \(Q\) 在曲线 \(x^{2} + (y + 2)^{2} = 1\) 上，那么 \(|PQ|\) 的最小值为（\quad）

\choices
    {\(\sqrt{5} - 1\)}
    {\(\frac{4}{\sqrt{5}} - 1\)}
    {\(2\sqrt{2} - 1\)}
    {\(\sqrt{2} - 1\)}
\end{problem}

\begin{answer}
A
\end{answer}

\begin{solution}
记圆的圆心为 \(O'=(0,-2)\). 由区域边界直线的交点可知，该区域是由
\[
y=2x+2,
\qquad y=\frac{x+1}{2},
\qquad y=2-x
\]
围成的三角形，其一个顶点为 \((-1,0)\). 对区域内任意点 \(P=(x,y)\)，由 \(x-2y+1\leqslant0\) 得
\[
x-2(y+2)=x-2y-4\leqslant-5.
\]
因此
\[
|PO'|=\sqrt{x^2+(y+2)^2}\geqslant\frac{|x-2y-4|}{\sqrt{5}}\geqslant\sqrt{5}.
\]
在 \(P=(-1,0)\) 处等号成立，所以点 \(O'\) 到该区域的最短距离为 \(\sqrt{5}\). 又 \(|QO'|=1\)，由三角不等式
\[
|PQ|\geqslant|PO'|-|QO'|\geqslant\sqrt{5}-1.
\]
从 \(O'\) 朝 \((-1,0)\) 方向取圆上的点 \(Q\)，即可使两次不等式同时取等号. 因此 \(|PQ|\) 的最小值为 \(\sqrt{5}-1\). 故选 A.
\end{solution}

\begin{problem}
半径为 1 的球面上的四点 \(A\)，\(B\)，\(C\)，\(D\) 是正四面体的顶点，则 \(A\) 与 \(B\) 两点间的球面距离为（\quad）

\choices
    {\(\arccos \left(-\frac{\sqrt{3}}{3}\right)\)}
    {\(\arccos \left(-\frac{\sqrt{6}}{3}\right)\)}
    {\(\arccos \left(-\frac{1}{3}\right)\)}
    {\(\arccos \left(-\frac{1}{4}\right)\)}
\end{problem}

\begin{answer}
C
\end{answer}

\begin{solution}
设正四面体的四个顶点相对于球心的单位位置向量分别为 \(\bs u_1\)，\(\bs u_2\)，\(\bs u_3\)，\(\bs u_4\). 正四面体的外接球球心也是其中心，所以四个位置向量的和为零. 由对称性，任意两个不同位置向量的内积相同，设为 \(t\). 于是
\[
0=\left|\bs u_1+\bs u_2+\bs u_3+\bs u_4\right|^2=4+2\times6t,
\]
从而 \(t=-\frac13\). 半径为 \(1\)，所以 \(A\)，\(B\) 与球心所成的圆心角 \(\theta\) 满足 \(\cos\theta=-\frac13\)，即 \(\theta=\arccos\left(-\frac13\right)\). 故选 C.
\end{solution}

\begin{problem}
如图，\(F_{1}\) 和 \(F_{2}\) 分别是双曲线 \( \frac{x^{2}}{a^{2}} - \frac{y^{2}}{b^{2}} = 1 \)（\(a > 0\)，\(b > 0\)）的两个焦点，\(A\) 和 \(B\) 是以 \(O\) 为圆心，以 \( |OF_{1}| \) 为半径的圆与该双曲线左支的两个交点，且 \( \triangle F_{2}AB \) 是等边三角形，则该双曲线的离心率为（\quad）

\FigureLayoutDeclare{img/2007/anhui_science/q9_fig1.png}{9.0cm}{22bp 0bp 18bp 15bp}

\bitmapfigure[max width=0.36\linewidth]{img/2007/anhui_science/q9_fig1.png}

\choices
    {\(\sqrt{3}\)}
    {\(\sqrt{5}\)}
    {\( \frac{\sqrt{5}}{2} \)}
    {\(1 + \sqrt{3}\)}
\end{problem}

\begin{answer}
D
\end{answer}

\begin{solution}
记双曲线的焦距参数为 \(c=\sqrt{a^2+b^2}\)，于是 \(F_1=(-c,0)\)，\(F_2=(c,0)\). 设左支上的交点为 \(A=(-x,y)\)，\(B=(-x,-y)\)，其中 \(x>0\)，\(y>0\). 由于 \(A\)，\(B\) 在半径为 \(c\) 的圆上，得
\[
x^2+y^2=c^2.
\]
又因为 \(\triangle F_2AB\) 为等边三角形，底边 \(AB=2y\)，且从 \(F_2\) 到底边中点的水平距离为 \(c+x\)，所以
\[
c+x=\sqrt{3}y.
\]
令 \(u=x/c\)，\(v=y/c\)，则 \(u^2+v^2=1\)，且 \(1+u=\sqrt{3}v\). 消去 \(v\) 得
\[
(1+u)^2=3(1-u^2),
\]
即 \(2u^2+u-1=0\). 由于 \(x>0\)，故 \(u=\frac12\)，从而 \(x=\frac c2\)，\(y=\frac{\sqrt{3}c}{2}\).

把点 \(A\) 代入双曲线方程，并令离心率 \(e=\frac ca\)，得
\[
\frac{c^2}{4a^2}-\frac{3c^2}{4b^2}=1,
\qquad b^2=c^2-a^2=a^2(e^2-1).
\]
因此
\[
\frac{e^2}{4}-\frac{3e^2}{4(e^2-1)}=1,
\]
化简得 \(e^4-8e^2+4=0\). 所以 \(e^2=4\pm2\sqrt{3}\). 双曲线的离心率满足 \(e>1\)，故
\[
e=\sqrt{4+2\sqrt{3}}=1+\sqrt{3}.
\]
故选 D.
\end{solution}

\begin{problem}
以 \(\Phi(x)\) 表示标准正态总体在区间 \((- \infty, x)\) 内取值的概率，若随机变量 \(\xi\) 服从正态分布 \(N(\mu, \sigma^2)\)，则概率 \(P(|\xi - \mu| < \sigma)\) 等于（\quad）

\choices
    {\(\Phi (\mu +\sigma) - \Phi (\mu -\sigma)\)}
    {\(\Phi (1) - \Phi (-1)\)}
    {\(\Phi \left(\frac{1 - \mu}{\sigma}\right)\)}
    {\(2\Phi (\mu +\sigma)\)}
\end{problem}

\begin{answer}
B
\end{answer}

\begin{solution}
令 \(Z=\frac{\xi-\mu}{\sigma}\)，则 \(Z\) 服从标准正态分布. 由于 \(\sigma>0\)，
\[
|\xi-\mu|<\sigma\iff-1<\frac{\xi-\mu}{\sigma}<1\iff-1<Z<1.
\]
按照 \(\Phi\) 的定义，所求概率为 \(\Phi(1)-\Phi(-1)\). 故选 B.
\end{solution}

\begin{problem}
定义在 \(\R\) 上的函数 \(f(x)\) 既是奇函数，又是周期函数，\(T\) 是它的一个正周期. 若将方程 \(f(x) = 0\) 在闭区间 \([-T, T]\) 上的根个数记为 \(n\)，则 \(n\) 可能为（\quad）

\choices
    {0}
    {1}
    {3}
    {5}
\end{problem}

\begin{answer}
D
\end{answer}

\begin{solution}
由奇函数性质，\(f(0)=0\). 再由周期性，\(f(-T)=f(0)=0\)，且 \(f(T)=f(0)=0\). 此外，\(T/2\) 与 \(-T/2\) 相差一个周期，所以
\[
f\left(\frac{T}{2}\right)=f\left(-\frac{T}{2}\right)=-f\left(\frac{T}{2}\right),
\]
从而 \(f(T/2)=f(-T/2)=0\). 因此区间 \([-T,T]\) 内至少有五个不同的根 \(-T\)，\(-T/2\)，\(0\)，\(T/2\)，\(T\). 例如取 \(f(x)=\sin\frac{2\pi x}{T}\)，它是奇函数且以 \(T\) 为正周期，并且在该区间内恰有这五个根. 因此 \(n\) 可能为 \(5\). 故选 D.
\end{solution}

\section{填空题}

\begin{problem}
若 \(\left(2x^{3} + \frac{1}{\sqrt{x}}\right)^{n}\) 的展开式中含有常数项，则最小的正整数 \(n\) 等于\fillinblank{}.
\end{problem}

\begin{answer}
7
\end{answer}

\begin{solution}
展开式通项中，若从第二项中取 \(k\) 个因子，则该项中 \(x\) 的幂为
\[
3(n-k)-\frac{k}{2}=3n-\frac{7k}{2}.
\]
含常数项的条件是 \(3n-\frac{7k}{2}=0\)，即 \(6n=7k\). 由于 \(n\)，\(k\) 为整数，且 \(0\leqslant k\leqslant n\)，最小的正整数解为 \(n=7\)，\(k=6\). 因此最小的正整数 \(n\) 等于 \(7\).
\end{solution}

\begin{problem}
在四面体 \(O - ABC\) 中，\(\overrightarrow{AB} = \bs{a}\)，\(\overrightarrow{OB} = \bs{b}\)，\(\overrightarrow{OC} = \bs{c}\)，\(D\) 为 \(BC\) 的中点，\(E\) 为 \(AD\) 的中点，则 \(\overrightarrow{OE} =\)\fillinblank{}（用 \(\bs{a}\)，\(\bs{b}\)，\(\bs{c}\) 表示）.
\end{problem}

\begin{answer}
\(\frac{-2\bs{a}+3\bs{b}+\bs{c}}{4}\)
\end{answer}

\begin{solution}
由 \(\overrightarrow{AB}=\bs{a}\) 和 \(\overrightarrow{OB}=\bs{b}\)，得
\[
\overrightarrow{OA}=\overrightarrow{OB}-\overrightarrow{AB}=\bs{b}-\bs{a}.
\]
因为 \(D\) 是 \(BC\) 的中点，所以
\[
\overrightarrow{OD}=\frac{\overrightarrow{OB}+\overrightarrow{OC}}{2}
=\frac{\bs{b}+\bs{c}}{2}.
\]
又因为 \(E\) 是 \(AD\) 的中点，所以
\[
\overrightarrow{OE}
=\frac{\overrightarrow{OA}+\overrightarrow{OD}}{2}
=\frac{\bs{b}-\bs{a}+\frac{\bs{b}+\bs{c}}{2}}{2}
=\frac{-2\bs{a}+3\bs{b}+\bs{c}}{4}.
\]
因此填 \(\frac{-2\bs{a}+3\bs{b}+\bs{c}}{4}\).
\end{solution}

\begin{problem}
如图，抛物线 \( y = -x^{2} + 1 \) 与 \( x \) 轴的正半轴交于点 \(A\)，将线段 \(OA\) 的 \( n \) 等分点从左至右依次记为 \(P_{1}\)，\(P_{2}\)，\(\cdots\)，\(P_{n-1}\)，过这些分点分别作 \( x \) 轴的垂线，与抛物线 \(C\) 的交点依次为 \(Q_{1}\)，\(Q_{2}\)，\(\cdots\)，\(Q_{n-1}\)，从而得到 \( n - 1 \) 个直角三角形 \(\triangle Q_{1}OP_{1}\)，\(\triangle Q_{2}P_{1}P_{2}\)，\(\cdots\)，\(\triangle Q_{n-1}P_{n-2}P_{n-1}\). 当 \( n \to \infty \) 时，这些三角形的面积之和的极限为\fillinblank{}.

\bitmapfigure[max width=0.48\linewidth]{img/2007/anhui_science/q14_fig1.png}
\end{problem}

\begin{answer}
\(\frac{1}{3}\)
\end{answer}

\begin{solution}
由抛物线 \(y=-x^2+1\) 与 \(x\) 轴正半轴的交点为 \(A=(1,0)\)，可知
\[
P_i=\left(\frac{i}{n},0\right),\qquad
Q_i=\left(\frac{i}{n},1-\frac{i^2}{n^2}\right)
\quad(1\leqslant i\leqslant n-1).
\]
第 \(i\) 个直角三角形的底边长为 \(\frac1n\)，高为 \(1-\frac{i^2}{n^2}\)，所以面积之和为
\[
S_n=\frac{1}{2n}\sum_{i=1}^{n-1}\left(1-\frac{i^2}{n^2}\right)
=\frac{n-1}{2n}-\frac{1}{2n^3}\sum_{i=1}^{n-1}i^2.
\]
利用 \(\sum_{i=1}^{n-1}i^2=\frac{(n-1)n(2n-1)}{6}\)，得
\[
\lim_{n\to\infty}S_n
=\lim_{n\to\infty}\left(\frac{n-1}{2n}
-\frac{(n-1)n(2n-1)}{12n^3}\right)
=\frac12-\frac16=\frac13.
\]
因此填 \(\frac13\).
\end{solution}

\begin{problem}
在正方体上任意选择 4 个顶点，它们可能是如下各种几何形体的 4 个顶点，这些几何形体是\fillinblank{}.（写出所有正确结论的编号）

\begin{circlelist}
\item 矩形；
\item 不是矩形的平行四边形；
\item 有三个面为等腰直角三角形，有一个面为等边三角形的四面体；
\item 每个面都是等边三角形的四面体；
\item 每个面都是直角三角形的四面体.
\end{circlelist}
\end{problem}

\begin{answer}
\(\circled{1}\circled{3}\circled{4}\circled{5}\)
\end{answer}

\begin{solution}
取正方体的顶点坐标为各坐标分量均为 \(0\) 或 \(1\) 的八个点.

取同一个面上的四个顶点，它们构成正方形，因而可以构成矩形，所以结论 \(\circled{1}\) 正确.

若四个顶点构成平行四边形，设从其中一个顶点出发的两条边向量为 \(\bs{u}\)，\(\bs{v}\)，则第四个顶点由 \(\bs{u}+\bs{v}\) 给出. 对每个坐标分量而言，为使这三个点仍都是正方体顶点，\(\bs{u}\) 与 \(\bs{v}\) 不可能在同一坐标分量上同时非零. 因此 \(\bs{u}\cdot\bs{v}=0\)，任意由正方体顶点构成的平行四边形都是矩形，不可能得到不是矩形的平行四边形，所以结论 \(\circled{2}\) 错误.

取四点 \((0,0,0)\)，\((1,0,0)\)，\((0,1,0)\)，\((0,0,1)\). 以第一个点为公共顶点的三个面都是等腰直角三角形，而另外一个面三个边长均为 \(\sqrt{2}\)，是等边三角形，所以结论 \(\circled{3}\) 正确.

取四点 \((0,0,0)\)，\((1,1,0)\)，\((1,0,1)\)，\((0,1,1)\). 任意两点间的距离都为 \(\sqrt{2}\)，所以它们构成每个面都是等边三角形的正四面体，结论 \(\circled{4}\) 正确.

取四点 \(O=(0,0,0)\)，\(U=(0,0,1)\)，\(V=(0,1,0)\)，\(W=(1,1,0)\). 三角形 \(OUV\) 和 \(OUW\) 分别在 \(O\) 点满足 \(\overrightarrow{OU}\cdot\overrightarrow{OV}=0\) 和 \(\overrightarrow{OU}\cdot\overrightarrow{OW}=0\)；三角形 \(OVW\) 在 \(V\) 点满足 \(\overrightarrow{VO}\cdot\overrightarrow{VW}=0\)；三角形 \(UVW\) 也在 \(V\) 点满足 \(\overrightarrow{VU}\cdot\overrightarrow{VW}=0\). 因此四个面都是直角三角形，结论 \(\circled{5}\) 正确.

所以所有正确结论的编号为 \(\circled{1}\circled{3}\circled{4}\circled{5}\).
\end{solution}

\section{解答题}

\begin{problem}
已知 \(0 < \alpha < \frac{\pi}{4}\)，\(\beta\) 为 \(f(x) = \cos \left(2x + \frac{\pi}{8}\right)\) 的最小正周期，\(\bs{a} = \left(\tan \left(\alpha + \frac{1}{4}\beta\right), -1\right)\)，\(\bs{b} = (\cos \alpha, 2)\)，且 \(\bs{a} \cdot \bs{b} = m\). 求 \(\frac{2\cos^2\alpha + \sin 2(\alpha + \beta)}{\cos \alpha - \sin \alpha}\) 的值.
\end{problem}

\begin{solution}
函数 \(f(x)=\cos\left(2x+\frac{\pi}{8}\right)\) 的最小正周期为
\[
\beta=\frac{2\pi}{2}=\pi.
\]
由点积条件，
\[
m=\bs{a}\cdot\bs{b}
=\tan\left(\alpha+\frac{\pi}{4}\right)\cos\alpha-2,
\]
所以 \(\tan\left(\alpha+\frac{\pi}{4}\right)\cos\alpha=m+2\).

又因为 \(0<\alpha<\frac{\pi}{4}\)，有 \(\cos\alpha-\sin\alpha>0\)，并且
\[
\tan\left(\alpha+\frac{\pi}{4}\right)
=\frac{\sin\alpha+\cos\alpha}{\cos\alpha-\sin\alpha}.
\]
由于 \(\beta=\pi\)，\(\sin2(\alpha+\beta)=\sin2\alpha\)，从而
\[
\frac{2\cos^2\alpha+\sin2(\alpha+\beta)}
{\cos\alpha-\sin\alpha}
=\frac{2\cos^2\alpha+2\sin\alpha\cos\alpha}
{\cos\alpha-\sin\alpha}
=2\tan\left(\alpha+\frac{\pi}{4}\right)\cos\alpha
=2(m+2)=2m+4.
\]
所以所求值为 \(2m+4\).
\end{solution}

\begin{problem}
如图，在六面体 \(ABCD - A_{1}B_{1}C_{1}D_{1}\) 中，四边形 \(ABCD\) 是边长为 2 的正方形，四边形 \(A_{1}B_{1}C_{1}D_{1}\) 是边长为 1 的正方形，\(DD_{1} \perp\) 平面 \(A_{1}B_{1}C_{1}D_{1}\)，\(DD_{1} \perp\) 平面 \(ABCD\)，\(DD_{1} = 2\).

\begin{enumerate}
\item 求证：\(A_{1}C_{1}\) 与 \(AC\) 共面，\(B_{1}D_{1}\) 与 \(BD\) 共面.
\item 求证：平面 \(A_{1}ACC_{1} \perp\) 平面 \(B_{1}BDD_{1}\).
\item 求二面角 \(A - BB_{1} - C\) 的大小.（用反三角函数值表示）
\end{enumerate}

\bitmapfigure[max width=0.38\linewidth]{img/2007/anhui_science/q17_fig1.png}
\end{problem}

\begin{solution}
以 \(D\) 为原点，以 \(DA\)，\(DC\)，\(DD_1\) 所在直线分别为 \(x\)，\(y\)，\(z\) 轴建立空间直角坐标系. 由题意可取
\[
D=(0,0,0),\quad A=(2,0,0),\quad C=(0,2,0),\quad B=(2,2,0),
\]
\[
D_1=(0,0,2),\quad A_1=(1,0,2),\quad C_1=(0,1,2),\quad B_1=(1,1,2).
\]
\begin{enumerate}
\item
\[
\overrightarrow{A_1C_1}=(-1,1,0)=\frac12\overrightarrow{AC},\qquad
\overrightarrow{B_1D_1}=(-1,-1,0)=\frac12\overrightarrow{BD}.
\]
所以 \(A_1C_1\parallel AC\)，\(B_1D_1\parallel BD\). 两组直线分别平行，故 \(A_1C_1\) 与 \(AC\) 共面，\(B_1D_1\) 与 \(BD\) 共面.

\item 平面 \(A_1ACC_1\) 的两个方向向量可取 \(\overrightarrow{AC}=(-2,2,0)\) 和 \(\overrightarrow{AA_1}=(-1,0,2)\)，所以其法向量可取
\[
\bs{n}_1=\overrightarrow{AC}\times\overrightarrow{AA_1}=(4,4,2).
\]
平面 \(B_1BDD_1\) 的两个方向向量可取 \(\overrightarrow{BD}=(-2,-2,0)\) 和 \(\overrightarrow{DD_1}=(0,0,2)\)，所以其法向量可取
\[
\bs{n}_2=\overrightarrow{BD}\times\overrightarrow{DD_1}=(-4,4,0).
\]
因为 \(\bs{n}_1\cdot\bs{n}_2=0\)，两平面的法向量垂直，所以平面 \(A_1ACC_1\perp\) 平面 \(B_1BDD_1\).

\item
设 \(M\) 为直线 \(BB_1\) 上满足 \(MA\perp BB_1\) 的点. 由于
\[
\overrightarrow{BA}\cdot\overrightarrow{BB_1}=2,
\qquad
|\overrightarrow{BB_1}|^2=6,
\]
可得
\[
\overrightarrow{BM}=\frac13\overrightarrow{BB_1}.
\]
又因为 \(\overrightarrow{BC}\cdot\overrightarrow{BB_1}=2\)，所以
\[
\overrightarrow{MC}\cdot\overrightarrow{BB_1}
=\overrightarrow{BC}\cdot\overrightarrow{BB_1}
-\frac13|\overrightarrow{BB_1}|^2
=2-\frac{6}{3}=0.
\]
故 \(MC\perp BB_1\)，从而 \(\angle AMC\) 是二面角 \(A-BB_1-C\) 的平面角. 由
\[
\overrightarrow{MA}=\left(\frac13,-\frac53,-\frac23\right),
\qquad
\overrightarrow{MC}=\left(-\frac53,\frac13,-\frac23\right),
\]
得
\[
\overrightarrow{MA}\cdot\overrightarrow{MC}=-\frac23,
\qquad
|\overrightarrow{MA}|^2=|\overrightarrow{MC}|^2=\frac{10}{3}.
\]
所以
\[
\cos\angle AMC=\frac{-\frac23}{\sqrt{\frac{10}{3}}\sqrt{\frac{10}{3}}}=-\frac15,
\]
从而 \(\theta=\arccos\left(-\frac15\right)\).
\end{enumerate}
\end{solution}

\begin{problem}
设 \(a \geqslant 0\)，\(f(x) = x - 1 - \ln^2 x + 2a \ln x\)（\(x > 0\)）.

\begin{enumerate}
\item 令 \(F(x) = xf'(x)\)，讨论 \(F(x)\) 在 \((0, +\infty)\) 内的单调性并求极值.
\item 求证：当 \(x > 1\) 时，恒有 \(x > \ln^{2} x - 2a \ln x + 1\).
\end{enumerate}
\end{problem}

\begin{solution}
\begin{enumerate}
\item
由
\[
f'(x)=1-\frac{2\ln x}{x}+\frac{2a}{x},
\]
得
\[
F(x)=xf'(x)=x-2\ln x+2a,
\qquad
F'(x)=1-\frac2x=\frac{x-2}{x}.
\]
当 \(0<x<2\) 时，\(F'(x)<0\)，所以 \(F(x)\) 在 \((0,2)\) 内单调递减；当 \(x>2\) 时，\(F'(x)>0\)，所以 \(F(x)\) 在 \((2,+\infty)\) 内单调递增. 因此 \(F(x)\) 在 \(x=2\) 处取得极小值
\[
F(2)=2-2\ln2+2a.
\]

\item 由第（1）问，\(F(x)\geqslant F(2)=2-2\ln2+2a>0\)，其中最后一个不等号由 \(a\geqslant0\) 和 \(\ln2<1\) 得到. 因此当 \(x>1\) 时，
\[
f'(x)=\frac{F(x)}{x}>0.
\]
所以 \(f(x)\) 在 \((1,+\infty)\) 内单调递增. 又 \(f(1)=0\)，故当 \(x>1\) 时 \(f(x)>f(1)=0\)，即
\[
x-1-\ln^2x+2a\ln x>0.
\]
移项即得 \(x>\ln^2x-2a\ln x+1\).
\end{enumerate}
\end{solution}

\begin{problem}
如图，曲线 \(G\) 的方程为 \(y^{2} = 2x\)（\(y \geqslant 0\)）. 以原点为圆心，以 \(t\)（\(t > 0\)）为半径的圆分别与曲线 \(G\) 和 \(y\) 轴的正半轴相交于点 \(A\) 与点 \(B\). 直线 \(AB\) 与 \(x\) 轴相交于点 \(C\).

\begin{enumerate}
\item 求点 \(A\) 的横坐标 \(a\) 与点 \(C\) 的横坐标 \(c\) 的关系式.
\item 设曲线 \(G\) 上点 \(D\) 的横坐标为 \(a + 2\)，求证：直线 \(CD\) 的斜率为定值.
\end{enumerate}

\bitmapfigure[max width=0.42\linewidth]{img/2007/anhui_science/q19_fig1.png}
\end{problem}

\begin{solution}
\begin{enumerate}
\item 设 \(A=(a,\sqrt{2a})\)，\(B=(0,t)\)，\(C=(c,0)\). 由 \(A\) 在圆 \(x^2+y^2=t^2\) 上，得
\[
a^2+2a=t^2,
\qquad
t=\sqrt{a(a+2)}.
\]
直线 \(BC\) 的截距式方程为
\[
\frac{x}{c}+\frac{y}{t}=1.
\]
由于点 \(A\) 在直线 \(BC\) 上，
\[
\frac{a}{c}+\frac{\sqrt{2a}}{t}=1,
\]
所以
\[
\frac{a}{c}=1-\sqrt{\frac{2}{a+2}}.
\]
令 \(s=\sqrt{a+2}\)，则 \(a=s^2-2=(s-\sqrt2)(s+\sqrt2)\)，从而
\[
c=\frac{a}{1-\frac{\sqrt2}{s}}
=\frac{(s-\sqrt2)(s+\sqrt2)}{\frac{s-\sqrt2}{s}}
=s(s+\sqrt2)
=a+2+\sqrt{2(a+2)}.
\]
因此 \(a\) 与 \(c\) 的关系式为
\[
c=a+2+\sqrt{2(a+2)}.
\]

\item 由 \(D\) 在曲线 \(G\) 上且横坐标为 \(a+2\)，得
\[
D=\left(a+2,\sqrt{2(a+2)}\right).
\]
由第（1）问，\(c=(a+2)+\sqrt{2(a+2)}\)，所以 \(D=(a+2,c-a-2)\). 于是直线 \(CD\) 的斜率为
\[
k_{CD}=\frac{c-a-2}{a+2-c}=-1.
\]
因此直线 \(CD\) 的斜率为定值 \(-1\).
\end{enumerate}
\end{solution}

\begin{problem}
在医学生物学试验中，经常以果蝇作为试验对象，一个装有 6 只果蝇的笼子里，不慎混入了两只苍蝇（此时笼内共有 8 只蝇子：6 只果蝇和 2 只苍蝇），只好把笼子打开一个小孔，让蝇子一只一只地往外飞. 直到两只苍蝇都飞出，再关闭小孔. 以 \(\xi\) 表示笼内还剩下的果蝇的只数.

\begin{enumerate}
\item 写出 \(\xi\) 的分布列（不要求写出计算过程）.
\item 求数学期望 \(E(\xi)\).
\item 求概率 \( P(\xi \geqslant E(\xi)) \).
\end{enumerate}
\end{problem}

\begin{solution}
\begin{enumerate}
\item 将两只苍蝇在八只蝇子飞出次序中的位置记为两个位置，任取两个位置的可能性相同. 若第二只苍蝇在第 \(k\) 个位置飞出，其中 \(2\leqslant k\leqslant8\)，则另一个苍蝇的位置可在前 \(k-1\) 个位置中任取一个，所以 \(P(\text{第二只苍蝇在第 }k\text{ 个位置飞出})=\frac{k-1}{\mathrm C_8^2}\). 此时剩余的果蝇数为 \(\xi=8-k\).

由 \(\xi=8-k\)，得到分布列
\begin{center}
\begin{tabular}{c|ccccccc}
\(\xi\) & \(0\) & \(1\) & \(2\) & \(3\) & \(4\) & \(5\) & \(6\) \\ \hline
\(P\) & \(\frac{7}{28}\) & \(\frac{6}{28}\) & \(\frac{5}{28}\) & \(\frac{4}{28}\) & \(\frac{3}{28}\) & \(\frac{2}{28}\) & \(\frac{1}{28}\)
\end{tabular}
\end{center}
\item
\[
E(\xi)=\frac{0\cdot7+1\cdot6+2\cdot5+3\cdot4+4\cdot3+5\cdot2+6\cdot1}{28}
=\frac{56}{28}=2.
\]

\item 因为 \(E(\xi)=2\)，所以
\[
P(\xi\geqslant E(\xi))
=P(\xi\geqslant2)
=\frac{5+4+3+2+1}{28}
=\frac{15}{28}.
\]
\end{enumerate}
\end{solution}

\begin{problem}
某国采用养老储备金制度. 公民在就业的第一年就交纳养老储备金，数目为 \(a_{1}\)；以后每年交纳的数目均比上一年增加 \(d\)（\(d > 0\)）. 因此，历年所交纳的储备金数目 \(a_{1}\)，\(a_{2}\)，\(\cdots\) 是一个公差为 \(d\) 的等差数列. 与此同时，国家给予优惠的计息政策，不仅采用固定利率，而且计算复利. 这就是说，如果固定年利率为 \(r\)（\(r > 0\)），那么，在第 \(n\) 年末，第一年所交纳的储备金就变为 \(a_{1}(1 + r)^{n-1}\)，第二年所交纳的储备金就变为 \(a_{2}(1 + r)^{n-2}\)，\(\cdots\)，以 \(T_{n}\) 表示到第 \(n\) 年末所累计的储备金总额.

\begin{enumerate}
\item 写出 \(T_{n}\) 与 \(T_{n-1}\)（\(n \geqslant 2\)）的递推关系式.
\item 求证：\(T_{n} = A_{n} + B_{n}\)，其中 \(\{A_{n}\}\) 是一个等比数列，\(\{B_{n}\}\) 是一个等差数列.
\end{enumerate}
\end{problem}

\begin{solution}
\begin{enumerate}
\item 由第 \(n-1\) 年末累计的储备金在一年后变为 \((1+r)T_{n-1}\)，再加上第 \(n\) 年交纳的 \(a_n=a_1+(n-1)d\)，得到
\[
T_n=(1+r)T_{n-1}+a_n=(1+r)T_{n-1}+a_1+(n-1)d,
\qquad n\geqslant2.
\]

\item 令
\[
A_n=\frac{(a_1r+d)(1+r)^n}{r^2},
\qquad
B_n=-\frac{d}{r}n-\frac{a_1}{r}-\frac{d}{r^2}.
\]
因为
\[
\frac{A_n}{A_{n-1}}=1+r,
\qquad
B_n-B_{n-1}=-\frac{d}{r},
\]
所以 \(\{A_n\}\) 是等比数列，\(\{B_n\}\) 是等差数列.

下面验证 \(T_n=A_n+B_n\). 由第（1）问的递推式及 \(T_1=a_1\)，先计算
\[
A_1+B_1
=\frac{(a_1r+d)(1+r)}{r^2}
-\frac{d}{r}-\frac{a_1}{r}-\frac{d}{r^2}
=a_1.
\]
又由 \(A_n=(1+r)A_{n-1}\) 和 \(B_n-B_{n-1}=-\frac{d}{r}\)，有
\[
(1+r)B_{n-1}+a_1+(n-1)d
-B_n
=rB_{n-1}+a_1+(n-1)d+\frac{d}{r}.
\]
代入 \(B_{n-1}=-\frac{d}{r}(n-1)-\frac{a_1}{r}-\frac{d}{r^2}\)，上式等于
\[
-d(n-1)-a_1-\frac{d}{r}+a_1+(n-1)d+\frac{d}{r}=0.
\]
因此 \(A_n+B_n=(1+r)(A_{n-1}+B_{n-1})+a_1+(n-1)d\)，且初值与 \(T_1\) 相同. 由递推关系的唯一性，\(T_n=A_n+B_n\) 对所有正整数 \(n\) 成立.
\end{enumerate}
\end{solution}
